\section{Conclusiones y Trabajo Futuro}

% [PENDIENTE: párrafo introductorio breve del capítulo una vez consolidados
%  los resultados de la corrida HALF/TOTAL.]

\subsubsection{Conclusiones}
% [PENDIENTE: redactar en tercera persona a partir de los resultados de la
%  Fase 3 (Capítulo V). Estructura sugerida alineada con la rúbrica
%  (relación resultados-metodología-objetivos):
%   1. Cumplimiento del Objetivo Específico 1 (arquitectura Transformer
%      adaptada al entrenamiento multi-serie con exógenas multimodales).
%   2. Cumplimiento del Objetivo Específico 2 (evaluación Cold-Start LOPO
%      global vs. baselines locales) y veredicto sobre la HIPÓTESIS del
%      Capítulo I.
%   3. Cumplimiento del Objetivo Específico 3 (integración multimodal:
%      exógenas dinámicas + estáticas geográficas).
%   4. Hallazgos metodológicos relevantes: rampas de puesta en marcha y
%      sensibilidad de la ventana de evaluación; efecto estacional
%      (dificultad invernal); integridad del target sin imputación.]

\subsubsection{Trabajos Futuros}
% [PENDIENTE: redactar fundamentado en los resultados. Candidatos ya
%  identificados durante el desarrollo:
%   - Ajuste fino (fine-tuning) del modelo global con los primeros días de
%     datos reales de la planta nueva (transfer learning en dos etapas).
%   - Incorporación de PatchTST cuando el ecosistema NeuralForecast
%     incorpore soporte de variables exógenas para dicha arquitectura.
%   - Extensión del prior regional de eficiencia con clustering climático
%     más fino que la Macrozona.
%   - Despliegue operacional (MLOps) explícitamente excluido del alcance.]
